% Root result promoted after two independent complete proof/source audits; original candidate preserved.
\subsection{The actual Cartier cubic and its rank strata}

\begin{lemma}[Identifying the Cartier cubic on a nonproduct surface]
\label{lem:root-nonordinary-Coble-identification-candidate}
\status{Proved under the stated nonproduct, nonordinary, characteristic-two hypotheses.}
Let $(A,\Theta)$ be a nonproduct principally polarized abelian surface
over an algebraically closed field of characteristic two, and suppose
that $A$ is nonordinary. Set $L=\mathcal O_A(3\Theta)$ and
$W=H^0(A,L)$. Use the typed theta basis and the alternating tensor
$T\in\bigwedge^3W^*$ of
Theorem~\ref{thm:ppav-level3-Cartier-symmetric}.
Let $M(v)$ be its contraction matrix and $q=q_\alpha$ the nonzero
Pfaffian cubic of Lemma~\ref{lem:DE-A3-theta-pfaffian}.
There is a $3\Theta^{(2)}$ embedding
\[
                 j:A^{(2)}\hookrightarrow\mathbf P(W)
\]
such that $q$ is, up to scalar, the Heisenberg-invariant Coble cubic
of this embedded surface. Its reduced singular locus is exactly
$j(A^{(2)})$. The rank-at-most-four locus of $M$ has this same
underlying set, and $M$ has rank exactly four at every point of it.
\end{lemma}
\begin{proof}
Write $F:A\to A^{(2)}$ for relative Frobenius and put
$E=F_*L$, a bundle of rank four. The established finite-duality
pairing is perfect and has global map
\[
 \beta:W\otimes W\longrightarrow W^{(2)}.
\]
The bundle $E$ is globally generated, as proved in
Proposition~\ref{prop:DE-A3-ppav-degree-one}. The map
$R:W^{(2)}\to W^*$ sends $e_a^{(2)}$ to $e_a^*$.
An evaluation functional of $L^{(2)}$ at $P\in A^{(2)}$, up to
its fibre trivialization, is an element of $(W^{(2)})^*$.
Define $j(P)$ by applying $(R^*)^{-1}$ to this functional.
Thus $j$ is precisely the complete $|L^{(2)}|$ embedding followed
by a fixed projective linear isomorphism; it does not identify
$A$ with $A^{(2)}$.

For any representative $v_P$ of $j(P)$, contraction $M(v_P)$ is
the perfect fibre pairing on $E_P$ pulled back by the surjection
$W\to E_P$. Its rank is therefore exactly four. In particular all
six-by-six Pfaffians vanish on the reduced embedded surface, so
the inclusion in the rank-at-most-four scheme holds scheme-theoretically.
On a chart $v_i\ne0$, the principal eight-by-eight Pfaffian is
$v_iq(v)$. Its differential is a linear combination of six-by-six
Pfaffians. At rank at most four both that Pfaffian and its
differential vanish; division by $v_i$ gives $q=dq=0$.
Hence $q$ is singular along $j(A^{(2)})$.

The finite Heisenberg action on the target of $\beta$ is the
coefficient-square representation, and $R$ identifies it with the
contragredient representation. Consequently the just-defined
embedding is equivariant for the same level-three Heisenberg group
acting on $\mathbf P(W)$. The cubic $q$ is invariant under this
group by the Pfaffian calculation already proved.

For clarity, the uniqueness used here is invariant uniqueness,
which is valid in characteristic two. Put $B=A^{(2)}$ and
$L_B=\mathcal O_B(3\Theta^{(2)})$. By
\citet[Theorem~2.4, p.~723]{Sekiguchi1977PNII}, applied with
the ample bundle $\mathcal O_B(\Theta^{(2)})$ and exponents $3,3$,
the multiplication map
\[
 \operatorname{Sym}^2H^0(B,L_B)\longrightarrow H^0(B,L_B^2)
\]
is surjective. Its dimensions are $45$ and $36$, respectively,
so $h^0(\mathcal I_{j(B)}(2))=9$.
\citet[Proposition~2.1, Remark~(3), Corollary~2.2, pp.~191--192]
{Beauville2003Coble} then gives a one-dimensional space of
Heisenberg-invariant cubics singular along $j(B)$.
The characteristic restriction there is only characteristic
different from three for this assertion. Its separate
Proposition~3.1 about unrestricted uniqueness is not needed.

The nonproduct principal surface is the Jacobian of its smooth
genus-two theta divisor, by
Lemma~\ref{lem:ppav-surface-dichotomy}. Choose a Weierstrass point
over the algebraically closed field. Translating the Weierstrass-point
representative of the principal theta divisor identifies its complete
$3\Theta$ embedding with the fixed complete embedding above; the induced
linear isomorphism transports the theta group and the tensor together.
In this identified embedding, the construction of
\citet[Section~4, Lemmas~4.1--4.4 and Proposition~4.5(a)]
{RainsSam2018Trivectors} gives a stable alternating tensor whose
rank-four surface is $j(B)$; the identification follows also from
its Propositions~5.1--5.2. Lemma~3.3 of that paper says its Pfaffian
cubic has singular locus $j(B)$ as a set. This cubic is invariant
under the same finite theta group. The preceding uniqueness makes
it a scalar multiple of $q$.

Every rank-at-most-four point of our matrix is singular on $q$ by
the local Pfaffian argument above. Conversely every point of
$j(B)$ has rank exactly four. Since the reduced singular locus
of $q$ is $j(B)$, this proves all the set-theoretic and exact-rank
claims. No equality of Jacobian schemes is inferred from the
quoted singular-locus statement.
\end{proof}

\begin{lemma}[A plane cubic rules out instability]
\label{lem:root-nonordinary-Cartier-stability-candidate}
\status{Proved under the stated nonproduct, nonordinary, characteristic-two hypotheses.}
The actual alternating Cartier tensor $T$ in the preceding lemma
is stable for $\operatorname{SL}(W^*)/\mu_3$.
Its rank-at-most-four scheme is smooth of dimension two and equals
$j(A^{(2)})$ scheme-theoretically. There are no nonzero points of
rank at most two, and $q=0$ is smooth at every rank-six point.
\end{lemma}
\begin{proof}
Suppose that $T$ is not stable. Apply
\citet[Proposition~2.12 and the proof of Lemma~3.2]
{RainsSam2018Trivectors}, with its $V_9=W^*$.
There is a six-dimensional subspace $U\subset W^*$ with
\[
 T\in\bigwedge^3U+\bigwedge^2U\otimes(W^*/U).
\]
On the plane $\Pi=\mathbf P(U^\perp)\subset\mathbf P(W)$,
the contraction matrix is supported on $U$. Its rank-at-most-four
locus is therefore cut out on $\Pi$ by one six-by-six Pfaffian,
a cubic polynomial. If that polynomial is identically zero, the
whole plane lies in the rank-at-most-four set, hence in $j(A^{(2)})$,
which is impossible: an abelian variety contains no rational curve.
Otherwise its zero scheme is a plane cubic supported in this
abelian surface.

Each irreducible component of the reduced support has degree at
most three. A line, a conic, or a singular integral plane cubic
has rational normalization, and would give a nonconstant map
from $\mathbf P^1$ to an abelian variety. Thus the only possibility
is a smooth integral plane cubic $E$, with multiplicity one; its
degree already exhausts the total degree three. After translating
an origin, its inclusion in the abelian surface is a homomorphism
from an elliptic curve, by the rigidity property of morphisms
between abelian varieties. Since the embedding is $3\Theta^{(2)}$,
\[
                         \Theta^{(2)}\cdot E=1.
\]
This forces the principal polarization to decompose, as follows.
If $i:E\hookrightarrow B=A^{(2)}$ is the inclusion, then
$i^*\lambda_\Theta i:E\to E^*$ is the degree-one polarization
of $E$ and is an isomorphism. Therefore
\[
 r=(i^*\lambda_\Theta i)^{-1}i^*\lambda_\Theta:B\to E
\]
is a retraction. The homomorphism
$E\times\ker r\to B$, $(e,b)\mapsto i(e)+b$, has the explicit
inverse $b\mapsto(r(b),b-i(r(b)))$ and is an isomorphism.
The kernel is consequently a connected elliptic curve. The two
factors are orthogonal for $\lambda_\Theta$ by construction, and
the degree-one principal polarization restricts principally to
both. This contradicts the nonproduct hypothesis.
Thus $T$ is stable.

Lemma~3.1 of the same source now gives smoothness and dimension two
for the actual rank-at-most-four scheme and excludes rank at most
two. The preceding lemma identifies its underlying set with the
closed smooth surface $j(B)$. Their scheme-theoretic inclusion
is an equality, since both schemes are reduced with that support.
Lemma~3.3 gives smoothness of $q=0$ away from that surface, as
claimed. This deduction uses smoothness of the Pfaffian scheme
only after stability has been established.
\end{proof}

\begin{corollary}[Only the rank-six Lang incidence remains]
\label{cor:root-nonordinary-Lang-rank-four-excluded-candidate}
\status{Proved under the stated nonproduct, nonordinary, characteristic-two hypotheses.}
In the affine coefficient model for the degree-one relative Lang
cover, write its nonzero fixed-vector equation as
\[
                    M(s)x^{[2]}=x,\qquad x\ne0.
\]
The part where $M(x^{[2]})$ has rank at most four cannot dominate
the nine-dimensional section-parameter space.
The part where that matrix has rank eight is irreducible and dominates,
with
an affine-line fibre over a nonempty open of the Fermat cone
$\sum x_a^3=0$.
Thus any further dominating component must come from the
rank-six compatibility locus. Global connectivity and full
monodromy are not asserted by this corollary.
\end{corollary}
\begin{proof}
Full symmetry of $T$ rewrites the equation as
$M(x^{[2]})s=x$. At rank four the allowed nonzero $x$ lie in
the Frobenius preimage of the affine cone over $j(A^{(2)})$.
This set has dimension three, since Frobenius is finite and
the projective surface has dimension two. For any such $x$,
the solution set in $s$ is empty or an affine space of dimension
$9-4=5$. The total incidence therefore has dimension at most
eight. Lower ranks do not occur for nonzero $x$ by the preceding
lemma. This proves nondominance of this entire part.

At rank eight, the kernel is $kx^{[2]}$, so compatibility is
exactly $x^t x^{[2]}=0$. The Fermat cubic in $\mathbf P^8$ is
smooth: its gradient is $(x_a^2)$ and has no projective zero.
It is integral, since distinct positive-degree hypersurface
components in projective space would meet and make their product
singular. Its nonzero affine cone is therefore irreducible of
dimension eight. The rank-eight open is nonempty because its
complement is the cubic $q(x^{[2]})=0$, and the proved
nonproportionality of $q$ with the Fermat cubic persists under
the coefficient-square-root twist: explicitly,
$q_\alpha(x^{[2]})=q_{\alpha^{[1/2]}}(x)^2$.
The affine-linear solutions in
$s$ form an affine-line torsor on this open, hence an irreducible
incidence of dimension nine.
The Jacobian of $M(s)x^{[2]}-x$ in the $x$ variables is the identity.
Thus its projection to the $s$ space is etale. On the nonempty
rank-eight open its image is open and hence dense, proving dominance.
The only other possible nonzero rank is six. Removing the
rank-at-most-four part does not bound that remaining incidence,
so the last sentence is a limitation, not a connectivity claim.
\end{proof}
